\documentclass[11pt,letterpaper]{article}
\usepackage[utf8]{inputenc}
\usepackage[T1]{fontenc}
\usepackage[spanish,es-noshorthands,es-tabla]{babel}
\usepackage[margin=2cm,top=2.5cm,bottom=2.5cm]{geometry}
\usepackage{amsmath,amssymb}
\usepackage{enumitem}
\usepackage{fancyhdr}
\usepackage{listings}
\usepackage{xcolor}
\usepackage{tcolorbox}
\usepackage{booktabs}
\usepackage{array}
\raggedbottom
\setlength{\headheight}{14pt}
\let\vspaceoriginal\vspace
\renewcommand{\vspace}[1]{\par\vspaceoriginal{#1}\par}
\pagestyle{fancy}
\fancyhf{}
\fancyhead[L]{Basic C Without Tears}
\fancyhead[C]{\NombreDocumento}
\fancyhead[R]{Stack}
\fancyfoot[L]{Basic-C-Without-Tears}
\fancyfoot[R]{\thepage}
\renewcommand{\headrulewidth}{0.4pt}
\renewcommand{\footrulewidth}{0.4pt}
\lstdefinestyle{customc}{
  language=C,
  numbers=left,
  stepnumber=1,
  numbersep=8pt,
  numberstyle=\tiny\color{gray},
  basicstyle=\footnotesize\ttfamily,
  keywordstyle=\bfseries\color{blue!80!black},
  commentstyle=\itshape\color{green!50!black},
  stringstyle=\color{red!70!black},
  breaklines=true,
  breakatwhitespace=true,
  tabsize=4,
  showstringspaces=false,
  frame=none
}
\newtcolorbox{solucionbox}{
  before=\par,
  after=\par,
  colback=white,
  colframe=black,
  arc=0mm,
  boxrule=0.6pt,
  left=3mm,
  right=3mm,
  top=2mm,
  bottom=2mm
}
\newif\ifpauta
\ifdefined\PAUTA\pautatrue\else\pautafalse\fi

\newcommand{\NombreDocumento}{Examen 04 -- Variante A}
\begin{document}
\begin{center}
    {\LARGE \textbf{Basic C Without Tears}} \\[0.2cm]
    {\Large \textbf{Examen 04 -- Variante A}} \\[0.12cm]
    \small Fundamentos de C, matrices, validación y algoritmos simples de strings \\[0.08cm]
    \textbf{Autor:} Stack
\end{center}
\vspace{0.2cm}
\begin{enumerate}[label=\textbf{\arabic*.}]
\item Un sensor guarda mediciones en la matriz:
\begin{lstlisting}[style=customc, numbers=none]
int S[2][3] = {{6, 2, 9}, {4, 7, 5}};
\end{lstlisting}
\begin{enumerate}[label=(\alph*)]
\item Indique el valor de \texttt{S[1][2]} y explique por qué \texttt{S[2][0]} no es un elemento válido.
\ifpauta\begin{solucionbox}\texttt{S[1][2]} vale 5. Los índices válidos de fila son 0 y 1; por eso \texttt{S[2][0]} excede los límites.\end{solucionbox}\else\vspace{1.5cm}\fi
\item Calcule el máximo de cada fila.
\ifpauta\begin{solucionbox}Fila 0: máximo 9; fila 1: máximo 7.\end{solucionbox}\else\vspace{2cm}\fi
\item Escriba código que encuentre e imprima el máximo de toda la matriz.
\ifpauta\begin{solucionbox}\begin{lstlisting}[style=customc]
int mayor = S[0][0];
for (int i = 0; i < 2; i++) {
    for (int j = 0; j < 3; j++) {
        if (S[i][j] > mayor) {
            mayor = S[i][j];
        }
    }
}
printf("Maximo: %d\n", mayor);
\end{lstlisting}El máximo global es 9. Se inicializa con un elemento real de la matriz.\end{solucionbox}\else\vspace{3cm}\fi
\end{enumerate}
\newpage
\item Para aceptar un paquete se conoce su peso \texttt{float peso} y se exige que sea mayor que 0 y menor o igual a 25 kg.
\begin{enumerate}[label=(\alph*)]
\item Escriba la condición que detecta un peso inválido.
\ifpauta\begin{solucionbox}\texttt{peso <= 0.0f || peso > 25.0f}.\end{solucionbox}\else\vspace{1.5cm}\fi
\item Escriba un fragmento que imprima si el paquete es aceptado o rechazado.
\ifpauta\begin{solucionbox}\begin{lstlisting}[style=customc]
if (peso > 0.0f && peso <= 25.0f) {
    printf("Paquete aceptado\n");
} else {
    printf("Paquete rechazado\n");
}
\end{lstlisting}\end{solucionbox}\else\vspace{2cm}\fi
\item Indique el resultado para \texttt{0.0}, \texttt{0.5}, \texttt{25.0} y \texttt{25.1}.
\ifpauta\begin{solucionbox}0.0: rechazado; 0.5: aceptado; 25.0: aceptado; 25.1: rechazado.\end{solucionbox}\else\vspace{2cm}\fi
\end{enumerate}
\newpage
\item Determine la salida del programa y describa qué representa el valor impreso.
\begin{lstlisting}[style=customc]
int A[2][2] = {{1, 3}, {5, 7}};
int suma = 0;
for (int i = 0; i < 2; i++) {
    suma += A[i][i];
}
printf("Suma diagonal: %d\n", suma);
\end{lstlisting}
\ifpauta\begin{solucionbox}El ciclo suma \texttt{A[0][0]} y \texttt{A[1][1]}: $1+7=8$.\par\textbf{Salida:} \texttt{Suma diagonal: 8}.\end{solucionbox}\else\vspace{3cm}\fi
\item Escriba una función que cuente la longitud de una cadena sin utilizar \texttt{strlen}. No cuente el terminador \texttt{'\\0'}.
\begin{center}\begin{tabular}{ll}\toprule\textbf{Entrada} & \textbf{Longitud}\\\midrule\texttt{"C"} & \texttt{1}\\\texttt{"hola"} & \texttt{4}\\\texttt{""} & \texttt{0}\\\bottomrule\end{tabular}\end{center}
\ifpauta\begin{solucionbox}\begin{lstlisting}[style=customc]
int longitud(const char texto[]) {
    int i = 0;
    while (texto[i] != '\0') {
        i++;
    }
    return i;
}
\end{lstlisting}Al detenerse justo en el terminador, el índice alcanzado coincide con la cantidad de caracteres.\end{solucionbox}\else\vspace{3cm}\fi
\end{enumerate}
\end{document}
